%\documentclass[letterpaper,12pt]{article}
\documentclass[preview]{standalone}

\usepackage{tabularx}
\usepackage{amsmath}
\usepackage{graphicx}
\usepackage[margin=1in,letterpaper]{geometry}
\usepackage{cite}
\usepackage[final]{hyperref}
\usepackage{datatool}
\usepackage{caption}
\usepackage{float}
\usepackage{gensymb}

\hypersetup{
        colorlinks=true,
        linkcolor=blue,
        citecolor=blue,
        filecolor=magenta,
        urlcolor=blue
}
\begin{document}

\begin{table}[h!]
 \centering
 \label{tab:fits}
 \renewcommand{\arraystretch}{1.2}
 \begin{tabular}{|c|c|c|c|c|c|}
 \hline
 Noyau & $A$ | V & $\lambda$ & $\omega$ | rad/s & T | s \\
 \hline
 Acier feuilleté     &
 $4.695 \pm 0.005$ &
 $111.1 \pm 0.2$ &
 $3376.4 \pm 0.2$ &
 $0.0018609 \pm 1 \cdot 10^{-7}$ \\

 Acier     &
 $4.35 \pm 0.01$ &
 $1179 \pm 5$ &
$4512 \pm 5$ &
 $0.0013925 \pm 1 \cdot 10^{-6}$ \\

 Air     &
 $4.618 \pm 0.005$ &
 $308.8 \pm 0.5$ &
 $7200.6 \pm 0.5$ &
 $0.0008726 \pm 6 \cdot 10^{-8}$ \\

 Aluminium     &
 $4.501 \pm 0.007$ &
 $910 \pm 2$ &
 $8894 \pm 2$ &
 $0.0007064 \pm 2 \cdot 10^{-7}$ \\
 \hline
 \end{tabular}
 \caption{Paramètres ajustés des oscillations avec le pont RLC et R = $10 \Omega$}
 \end{table}


\end{document}

